\section{RQ1.5 Azure Kubernetes Service Transfer Validation}
\label{sec:rq15}

\subsection{Research Question}

\begin{quote}
\textit{Does the local Kubernetes latency pattern observed for the coarse-stage
ResNet-18 microservice chain transfer to Azure Kubernetes Service, and how do the
absolute and relative costs change under managed cloud execution?}
\end{quote}

RQ1.5 moves the predefined RQ1.4 chain family from local \texttt{kind} Kubernetes to
Azure Kubernetes Service (AKS). The purpose is not to select a new decomposition, but to
test whether the architectural ordering observed in local Kubernetes remains credible
when the same synchronously chained inference workload is executed on managed cloud
infrastructure. The stage therefore compares absolute latency, relative overhead, and
condition ordering between the frozen RQ1.4 local Kubernetes run and the frozen RQ1.5
AKS run.

The literature framing for transfer-validation between local Kubernetes and
managed cloud execution is discussed in \cref{sec:rw-cloud-native}.

\subsection{Experimental Setup}

The thesis-facing RQ1.5 benchmark used the frozen controlled AKS
transfer-validation benchmark exported on 2026-04-22. The environment snapshot records a Linux Azure runtime
(\texttt{Linux-5.15.0-1102-azure-x86\_64-with-glibc2.41}), Python~3.10.20, and
PyTorch~2.11.0+cu130. The deployment ran in namespace \texttt{rq15} on a single-node
AKS cluster using Kubernetes~1.34, a \texttt{Standard\_D8s\_v3} node in the
\texttt{swedencentral} Azure region, and \texttt{kubenet} networking. Deployment
metadata confirms that pod colocation was enforced: for every condition, the benchmark
client and all service pods were scheduled on the same AKS node,
\texttt{aks-rq15pool-24384142-vmss000000}. All AKS conditions used the same pinned
container image digest:
\texttt{sha256:df3cf16f8760493993f8}\allowbreak
\texttt{ad5ba103bc5508a130180baf5130}\allowbreak
\texttt{ec44542a44513a8a}.

\begin{table}[ht]
\centering
\caption{RQ1.5 AKS benchmark configuration}
\label{tab:rq15_config}
\begin{tabular}{p{0.34\linewidth}p{0.58\linewidth}}
\toprule
\textbf{Setting} & \textbf{Value} \\
\midrule
Runtime & Azure Kubernetes Service benchmark \\
Cluster & Single-node AKS cluster, Kubernetes~1.34 \\
Node type & \texttt{Standard\_D8s\_v3} \\
Region & \texttt{swedencentral} \\
Network plugin & \texttt{kubenet} \\
Namespace & \texttt{rq15} \\
Conditions & Five predefined configurations from \texttt{monolithic\_k8s\_1svc} to \texttt{chain\_5svc} \\
Execution mode & In-cluster rolling orchestration with teardown between conditions \\
Model & ResNet-18 (pretrained) \\
Device & CPU \\
Precision & FP32 \\
Input shape & $1 \times 3 \times 224 \times 224$ \\
Rounds & 5 \\
Warmup iterations & 50 \\
Measured iterations & 200 \\
Cooldown between conditions & 5 seconds \\
Random seed & 42 \\
Communication & gRPC (in-cluster Kubernetes) \\
gRPC port & 50051 \\
Maximum message size & 16~MB \\
Service pod CPU request / limit & 1 core / 1 core \\
Service pod memory request / limit & 1~GiB / 1~GiB \\
Client CPU request / limit & 500m / 500m \\
Client memory request / limit & 1~GiB / 1~GiB \\
Pod QoS & \texttt{Guaranteed} for service pods \\
Image pull policy & \texttt{Always} \\
Carry-forward & Not applicable; RQ1.5 evaluates a predefined chain family \\
Parity validation & Local segment validation and live gRPC chain validation; \texttt{num\_inputs} $= 5$, \texttt{atol} $= 1 \times 10^{-5}$ \\
Warmup calibration & Trailing window $= 10$, CV threshold $= 0.02$ \\
\bottomrule
\end{tabular}
\end{table}

CPU stabilisation controls were requested consistently with the preceding stages, and
were applied where the managed AKS runtime allowed. Threading controls for PyTorch,
OpenMP, MKL, and OpenBLAS were requested and observed at one thread. CPU affinity was
applied to one physical core with SMT excluded; the observed affinity set was core~0.
Process priority \texttt{nice~-5} was requested but could not be applied due to
insufficient privileges, so the benchmark ran at \texttt{nice=0}. CPU governor and
turbo controls were unavailable through the exposed sysfs interface in AKS. These
constraints are treated as realistic cloud-environment limitations of managed or
virtualised execution, not as architectural effects of the microservice decomposition.

Before timing, both local chain-segment validation and live gRPC validation passed for
every condition with \texttt{max\_abs\_diff} $= 0.0$ and
\texttt{atol} $= 1 \times 10^{-5}$. This supports functional equivalence across the
monolithic and chained AKS deployments.

\subsection{Conditions}

RQ1.5 reused the same predefined chain family as RQ1.4:

\begin{itemize}
    \item \textbf{\texttt{monolithic\_k8s\_1svc}}: Full ResNet-18 inference in a single Kubernetes service.
    \item \textbf{\texttt{chain\_2svc}}: Two services with a split after \texttt{layer2}. Cumulative transferred activation per request: 980~KB.
    \item \textbf{\texttt{chain\_3svc}}: Three services with splits after \texttt{layer1} and \texttt{layer3}. Cumulative transferred activation per request: 1568~KB.
    \item \textbf{\texttt{chain\_4svc}}: Four services with splits after \texttt{layer1}, \texttt{layer2}, and \texttt{layer3}. Cumulative transferred activation per request: 1960~KB.
    \item \textbf{\texttt{chain\_5svc}}: Five services with splits after \texttt{layer1}, \texttt{layer2}, \texttt{layer3}, and \texttt{layer4}. Cumulative transferred activation per request: 2058~KB.
\end{itemize}

Carry-forward is not applicable in this stage. The experiment evaluates a fixed
deployment family in AKS rather than choosing a new split point.

\subsection{Results}

\begin{table}[ht]
\centering
\caption{RQ1.5 AKS summary of mean latency and overhead vs.\ AKS monolithic baseline}
\label{tab:rq15_results}
\resizebox{\linewidth}{!}{%
\begin{tabular}{lccccc}
\toprule
\textbf{Condition} & \textbf{Mean (ms)} & \textbf{Median (ms)} & \textbf{Std (ms)} & \textbf{p95 (ms)} & \textbf{Overhead vs.\ baseline} \\
\midrule
\texttt{monolithic\_k8s\_1svc} & 65.619 & 65.312 & 1.291 & 67.693 & --- \\
\texttt{chain\_2svc}           & 68.900 & 68.599 & 1.578 & 71.269 & +3.280 / +5.0\% \\
\texttt{chain\_3svc}           & 71.884 & 71.488 & 1.701 & 74.798 & +6.264 / +9.5\% \\
\texttt{chain\_4svc}           & 74.172 & 73.823 & 1.622 & 76.712 & +8.553 / +13.0\% \\
\texttt{chain\_5svc}           & 75.695 & 75.339 & 1.826 & 77.993 & +10.075 / +15.4\% \\
\bottomrule
\end{tabular}
}
\end{table}

\begin{table}[ht]
\centering
\caption{RQ1.5 AKS overhead and inferred non-compute/platform cost breakdown}
\label{tab:rq15_overhead}
\resizebox{\linewidth}{!}{%
\begin{tabular}{lcccc}
\toprule
\textbf{Condition} & \textbf{Overhead (ms)} & \textbf{Overhead (\%)} & \textbf{Cumul.\ Activation (KB)} & \textbf{Inferred Non-Compute / Platform (ms)} \\
\midrule
\texttt{chain\_2svc} & 3.280  & 5.0  & 980  & 4.053 \\
\texttt{chain\_3svc} & 6.264  & 9.5  & 1568 & 6.337 \\
\texttt{chain\_4svc} & 8.553  & 13.0 & 1960 & 8.041 \\
\texttt{chain\_5svc} & 10.075 & 15.4 & 2058 & 9.381 \\
\bottomrule
\end{tabular}
}
\end{table}

\begin{table}[ht]
\centering
\caption{RQ1.5 transfer comparison between local Kubernetes and AKS}
\label{tab:rq15_transfer}
\resizebox{\linewidth}{!}{%
\begin{tabular}{lcccc}
\toprule
\textbf{Condition} & \textbf{Local mean (ms)} & \textbf{AKS mean (ms)} & \textbf{Local overhead} & \textbf{AKS overhead} \\
\midrule
\texttt{monolithic\_k8s\_1svc} & 81.447 & 65.619 & 0.0\%  & 0.0\% \\
\texttt{chain\_2svc}           & 84.189 & 68.900 & 3.4\%  & 5.0\% \\
\texttt{chain\_3svc}           & 87.608 & 71.884 & 7.6\%  & 9.5\% \\
\texttt{chain\_4svc}           & 92.312 & 74.172 & 13.3\% & 13.0\% \\
\texttt{chain\_5svc}           & 92.624 & 75.695 & 13.7\% & 15.4\% \\
\bottomrule
\end{tabular}
}
\end{table}

\begin{figure}[htbp]
\centering
\resizebox{0.92\textwidth}{!}{%
\begin{tikzpicture}[
    font=\normalsize,
    axis/.style={thick, draw=black!70, -{Latex[length=2mm]}},
    grid/.style={draw=black!12, thin},
    localpoint/.style={
        circle,
        draw=blue!65!black,
        fill=blue!18,
        thick,
        minimum size=5.5pt,
        inner sep=0pt
    },
    akspoint/.style={
        rectangle,
        draw=green!45!black,
        fill=green!18,
        thick,
        minimum size=5.5pt,
        inner sep=0pt
    },
    localseries/.style={very thick, draw=blue!65!black},
    aksseries/.style={very thick, draw=green!45!black},
    pointlabel/.style={font=\small\bfseries, align=center},
    xlabel/.style={font=\footnotesize, text=black!75, align=center},
    ylabel/.style={font=\normalsize, text=black!75},
    legendlabel/.style={font=\footnotesize, text=black!75}
]

% y scale: y = overhead percentage * 0.25
% Local Kubernetes overheads (Stage K frozen): 0.0, 2.9,  8.2,  11.7, 14.1
% AKS overheads        (Stage C  frozen AKS):   0.0, 3.7,  6.5,  10.6, 13.7

\foreach \pct/\y in {0/0,5/1.25,10/2.5,15/3.75} {
    \draw[grid] (0,\y) -- (10.4,\y);
    \node[anchor=east, font=\footnotesize, text=black!65] at (-0.18,\y) {\pct};
}

\draw[axis] (0,0) -- (10.75,0);
\draw[axis] (0,0) -- (0,4.15);

\node[ylabel, rotate=90] at (-1.05,2.05) {Overhead vs.\ environment baseline (\%)};
\node[ylabel] at (5.2,-1.25) {Condition};

\coordinate (l1) at (1,0.000);
\coordinate (l2) at (3,0.725);
\coordinate (l3) at (5,2.050);
\coordinate (l4) at (7,2.925);
\coordinate (l5) at (9,3.525);

\coordinate (a1) at (1,0.000);
\coordinate (a2) at (3,0.925);
\coordinate (a3) at (5,1.625);
\coordinate (a4) at (7,2.650);
\coordinate (a5) at (9,3.425);

\draw[localseries] (l1) -- (l2) -- (l3) -- (l4) -- (l5);
\draw[aksseries] (a1) -- (a2) -- (a3) -- (a4) -- (a5);

\node[localpoint] at (l1) {};
\node[localpoint] at (l2) {};
\node[localpoint] at (l3) {};
\node[localpoint] at (l4) {};
\node[localpoint] at (l5) {};

\node[akspoint] at (a1) {};
\node[akspoint] at (a2) {};
\node[akspoint] at (a3) {};
\node[akspoint] at (a4) {};
\node[akspoint] at (a5) {};

\node[pointlabel, text=blue!65!black, below right=0.06cm and 0.05cm of l2] {2.9};
\node[pointlabel, text=blue!65!black, above=0.10cm of l3] {8.2};
\node[pointlabel, text=blue!65!black, above=0.10cm of l4] {11.7};
\node[pointlabel, text=blue!65!black, above=0.10cm of l5] {14.1};

\node[pointlabel, text=green!45!black, above=0.10cm of a2] {3.7};
\node[pointlabel, text=green!45!black, below right=0.06cm and 0.05cm of a3] {6.5};
\node[pointlabel, text=green!45!black, below=0.16cm of a4] {10.6};
\node[pointlabel, text=green!45!black, below=0.16cm of a5] {13.7};

\node[xlabel] at (1,-0.45) {\texttt{1svc}\\mono};
\node[xlabel] at (3,-0.45) {\texttt{2svc}};
\node[xlabel] at (5,-0.45) {\texttt{3svc}};
\node[xlabel] at (7,-0.45) {\texttt{4svc}};
\node[xlabel] at (9,-0.45) {\texttt{5svc}};

\draw[localseries] (6.25,0.70) -- (6.95,0.70);
\node[localpoint] at (6.60,0.70) {};
\node[legendlabel, anchor=west] at (7.15,0.70) {Local Kubernetes};

\draw[aksseries] (6.25,0.28) -- (6.95,0.28);
\node[akspoint] at (6.60,0.28) {};
\node[legendlabel, anchor=west] at (7.15,0.28) {AKS};

\end{tikzpicture}

}
\caption{RQ1.5 normalised overhead transfer comparison between local Kubernetes and AKS. Overheads are computed relative to each environment's own monolithic Kubernetes baseline, so the figure visualises directional transfer of the service-count trend rather than raw latency equivalence.}
\label{fig:rq15_transfer_overhead}
\end{figure}

\Cref{fig:rq15_transfer_overhead} visualises the transfer comparison in
\cref{tab:rq15_transfer}. The overhead trend is similar across environments even though
the absolute millisecond values differ.

The rank ordering is identical between the frozen local Kubernetes and AKS runs:
\texttt{monolithic\_k8s\_1svc}, \texttt{chain\_2svc}, \texttt{chain\_3svc},
\texttt{chain\_4svc}, and \texttt{chain\_5svc}. Absolute latency is lower in the AKS
run for every condition, but the within-environment ordering and the direction of the
overhead trend are preserved.

\begin{table}[ht]
\centering
\caption{RQ1.5 AKS cross-round consistency}
\label{tab:rq15_rounds}
\resizebox{\linewidth}{!}{%
\begin{tabular}{lcccc}
\toprule
\textbf{Condition} & \textbf{Grand Mean (ms)} & \textbf{Round Std (ms)} & \textbf{Round Range (ms)} & \textbf{Round CV} \\
\midrule
\texttt{monolithic\_k8s\_1svc} & 65.619 & 0.1714 & 0.4432 & 0.0026 \\
\texttt{chain\_2svc}           & 68.900 & 0.1628 & 0.3750 & 0.0024 \\
\texttt{chain\_3svc}           & 71.884 & 0.1290 & 0.3199 & 0.0018 \\
\texttt{chain\_4svc}           & 74.172 & 0.2225 & 0.5384 & 0.0030 \\
\texttt{chain\_5svc}           & 75.695 & 0.1291 & 0.3284 & 0.0017 \\
\bottomrule
\end{tabular}
}
\end{table}

The AKS run was highly stable across measured rounds. Round CVs ranged from 0.0017 to
0.0030, which is lower than the corresponding local Kubernetes CVs reported in RQ1.4.
Warmup calibration indicated that each run reached a stable window at least once before
measurement, while the cross-round CVs of the measured iterations provide the stronger
stability evidence for the final reported means. As in the previous stage, effect sizes
and Mann-Whitney tests were computed from pooled per-iteration data and are treated as
supplementary because $N = 1{,}000$ per condition makes very small differences
statistically significant. Cross-round stability, parity validation, and preserved
ordering are the primary evidence for this stage.

The interpretation of these results and the answer to RQ1.5 are presented in
\cref{sec:analysis-rq15}.
